Santos, \textit{A Terra Dá, A Terra Quer}.
Na p. 36: \enquote{quando cheguei ao território em que estou hoje, já existiam outros compartilhantes que nos recepcionaram. Na Caatinga, os umbuzeiros nos recepcionaram. Eles compartilharam seus frutos, suas folhas e suas raízes quando chegamos, e não trouxemos nada para os umbuzeiros. Eles já eram nativos daqui, viemos habitar esta terra depois deles. Foi assim com os pássaros, foi assim com uma planta chamada pinhão -- que não é o pinhão manso, é um pinhão cuidado por nós, ditos humanos, que as juritis adoram.  Elas comem esses pinhões e, vez por outra, pegamos uma juriti. O pinhão compartilha com a juriti, a juriti compartilha conosco, e nós vamos compartilhar de novo com o pinhão. Agora que já estamos aqui há mais tempo, entramos também no ciclo local de compartilhamento.}